
\documentclass{article}
\usepackage{colortbl}
\usepackage{makecell}
\usepackage{multirow}
\usepackage{supertabular}

\begin{document}

\newcounter{utterance}

\centering \large Interaction Transcript for game `cladder', experiment `full\_v1.5\_default', episode 4371 with refineonly{-}Qwen3.5{-}9B.
\vspace{24pt}

{ \footnotesize  \setcounter{utterance}{1}
\setlength{\tabcolsep}{0pt}
\begin{supertabular}{c@{$\;$}|p{.15\linewidth}@{}p{.15\linewidth}p{.15\linewidth}p{.15\linewidth}p{.15\linewidth}p{.15\linewidth}}
    \# & \multicolumn{2}{c}{Player} && \multicolumn{2}{c}{Game Master} \\
    \hline

    \theutterance \stepcounter{utterance}  
    & & & \multicolumn{4}{p{0.6\linewidth}}{
        \cellcolor[rgb]{0.9,0.9,0.9}{
            \makecell[{{p{\linewidth}}}]{
                \texttt{\tiny{[P1$\langle$GM]}}
                \texttt{You are an expert in causal inference. The following question is not a typical commonsense query, but rather a meticulously designed question created by a professor specializing in causal inference, intended to assess the students' mastery of the course content.} \\
\texttt{Imagine a self{-}contained, hypothetical world with only the following conditions, and without any unmentioned factors or causal relationships: Gender has a direct effect on smoking and peanut allergy. Smoking has a direct effect on tar deposit. Tar deposit has a direct effect on peanut allergy. Gender is unobserved. For nonsmokers, the probability of high tar deposit is 28\%. For smokers, the probability of high tar deposit is 58\%. For nonsmokers and with no tar deposit, the probability of being allergic to peanuts is 49\%. For nonsmokers and with high tar deposit, the probability of being allergic to peanuts is 77\%. For smokers and with no tar deposit, the probability of being allergic to peanuts is 27\%. For smokers and with high tar deposit, the probability of being allergic to peanuts is 62\%. The overall probability of smoking is 58\%. Will smoking increase the chance of being allergic to peanuts?} \\
\texttt{Start your answer with 'Yes' or 'No', followed by additional reasoning or evidence to support your explanation.} \\
            }
        }
    }
     \\ \\

    \theutterance \stepcounter{utterance}  
    & & \multicolumn{4}{p{0.6\linewidth}}{
        \cellcolor[rgb]{0.95,0.95,0.95}{
            \makecell[{{p{\linewidth}}}]{
                \texttt{\tiny{[P1$\leftrightarrow$P]}}
                \texttt{{[}DRAFT{]}} \\
\texttt{No} \\
\\ 
\texttt{**Reasoning:**} \\
\\ 
\texttt{To determine if smoking increases the chance of being allergic to peanuts, we must calculate the causal effect of smoking on peanut allergy. This requires comparing the probability of peanut allergy under the intervention of smoking (\$P(A|do(S=1))\$) versus non{-}smoking (\$P(A|do(S=0))\$).} \\
\\ 
\texttt{In this causal graph, **Gender** is a confounder because it has direct arrows to both Smoking (\$S\$) and Peanut Allergy (\$A\$). Additionally, there is a path from Smoking to Tar Deposit to Peanut Allergy. Since Gender is unobserved, we cannot simply average the conditional probabilities provided; we must account for the distribution of Gender in the population to perform the adjustment.} \\
\\ 
\texttt{Let \$G\$ represent Gender, \$S\$ represent Smoking, \$T\$ represent Tar Deposit, and \$A\$ represent Peanut Allergy.} \\
\texttt{The causal effect is defined as:} \\
\texttt{\$\$P(A|do(S=1)) {-} P(A|do(S=0))\$\$} \\
\\ 
\texttt{Using the backdoor adjustment criterion, we adjust for the confounder \$G\$:} \\
\texttt{\$\$P(A|do(S=s)) = \textbackslash{}sum\_g P(A|S=s, G=g)P(G=g)\$\$} \\
\\ 
\texttt{We are given the following conditional probabilities for Peanut Allergy (\$A\$):} \\
\texttt{1.  \$P(A|S=0, T=0) = 0.49\$} \\
\texttt{2.  \$P(A|S=0, T=1) = 0.77\$} \\
\texttt{3.  \$P(A|S=1, T=0) = 0.27\$} \\
\texttt{4.  \$P(A|S=1, T=1) = 0.62\$} \\
\\ 
\texttt{We are also given the probabilities for Tar Deposit (\$T\$) given Smoking (\$S\$):} \\
\texttt{1.  \$P(T=1|S=0) = 0.28 \textbackslash{}implies P(T=0|S=0) = 0.72\$} \\
\texttt{2.  \$P(T=1|S=1) = 0.58 \textbackslash{}implies P(T=0|S=1) = 0.42\$} \\
\\ 
\texttt{First, let's calculate the probability of Peanut Allergy for a specific gender group \$g\$ if they were to smoke (\$do(S=1)\$) or not smoke (\$do(S=0)\$).} \\
\texttt{Let \$p\_g = P(A|G=g, S=0)\$ be the baseline probability of allergy for gender \$g\$ if they don't smoke.} \\
\texttt{However, the problem states "Gender has a direct effect on smoking and peanut allergy." It does not explicitly give the probability of allergy *given only gender*. Instead, it gives the probability of allergy given gender, smoking, and tar deposit.} \\
\\ 
\texttt{Let's re{-}evaluate the structure. The question asks if smoking *incre} \\
            }
        }
    }
    & \\ \\

    \theutterance \stepcounter{utterance}  
    & \multicolumn{4}{p{0.6\linewidth}}{
        \cellcolor[rgb]{0.561,0.737,0.561}{
            \makecell[{{p{\linewidth}}}]{
                \texttt{\tiny{[P1$\rangle$GM]}}
                \texttt{No} \\
\\ 
\texttt{**Reasoning:**} \\
\\ 
\texttt{To determine if smoking increases the chance of being allergic to peanuts, we must calculate the causal effect of smoking on peanut allergy. This requires comparing the probability of peanut allergy under the intervention of smoking (\$P(A|do(S=1))\$) versus non{-}smoking (\$P(A|do(S=0))\$).} \\
\\ 
\texttt{In this causal graph, **Gender** is a confounder because it has direct arrows to both Smoking (\$S\$) and Peanut Allergy (\$A\$). Additionally, there is a path from Smoking to Tar Deposit to Peanut Allergy. Since Gender is unobserved, we cannot simply average the conditional probabilities provided; we must account for the distribution of Gender in the population to perform the adjustment.} \\
\\ 
\texttt{Let \$G\$ represent Gender, \$S\$ represent Smoking, \$T\$ represent Tar Deposit, and \$A\$ represent Peanut Allergy.} \\
\texttt{The causal effect is defined as:} \\
\texttt{\$\$P(A|do(S=1)) {-} P(A|do(S=0))\$\$} \\
\\ 
\texttt{Using the backdoor adjustment criterion, we adjust for the confounder \$G\$:} \\
\texttt{\$\$P(A|do(S=s)) = \textbackslash{}sum\_g P(A|S=s, G=g)P(G=g)\$\$} \\
\\ 
\texttt{We are given the following conditional probabilities for Peanut Allergy (\$A\$):} \\
\texttt{1.  \$P(A|S=0, T=0) = 0.49\$} \\
\texttt{2.  \$P(A|S=0, T=1) = 0.77\$} \\
\texttt{3.  \$P(A|S=1, T=0) = 0.27\$} \\
\texttt{4.  \$P(A|S=1, T=1) = 0.62\$} \\
\\ 
\texttt{We are also given the probabilities for Tar Deposit (\$T\$) given Smoking (\$S\$):} \\
\texttt{1.  \$P(T=1|S=0) = 0.28 \textbackslash{}implies P(T=0|S=0) = 0.72\$} \\
\texttt{2.  \$P(T=1|S=1) = 0.58 \textbackslash{}implies P(T=0|S=1) = 0.42\$} \\
\\ 
\texttt{First, let's calculate the probability of Peanut Allergy for a specific gender group \$g\$ if they were to smoke (\$do(S=1)\$) or not smoke (\$do(S=0)\$).} \\
\texttt{Let \$p\_g = P(A|G=g, S=0)\$ be the baseline probability of allergy for gender \$g\$ if they don't smoke.} \\
\texttt{However, the problem states "Gender has a direct effect on smoking and peanut allergy." It does not explicitly give the probability of allergy *given only gender*. Instead, it gives the probability of allergy given gender, smoking, and tar deposit.} \\
\\ 
\texttt{Let's re{-}evaluate the structure. The question asks if smoking *incre} \\
            }
        }
    }
    & & \\ \\

    \theutterance \stepcounter{utterance}  
    & & & \multicolumn{2}{p{0.3\linewidth}}{
        \cellcolor[rgb]{0.9,0.9,0.9}{
            \makecell[{{p{\linewidth}}}]{
                \texttt{\tiny{[GM$|$GM]}}
                \texttt{No} \\
            }
        }
    }
    & & \\ \\

    \theutterance \stepcounter{utterance}  
    & & & \multicolumn{2}{p{0.3\linewidth}}{
        \cellcolor[rgb]{0.9,0.9,0.9}{
            \makecell[{{p{\linewidth}}}]{
                \texttt{\tiny{[GM$|$GM]}}
                \texttt{yes} \\
            }
        }
    }
    & & \\ \\

    \theutterance \stepcounter{utterance}  
    & & & \multicolumn{2}{p{0.3\linewidth}}{
        \cellcolor[rgb]{0.9,0.9,0.9}{
            \makecell[{{p{\linewidth}}}]{
                \texttt{\tiny{[GM$|$GM]}}
                \texttt{game\_result = LOSE} \\
            }
        }
    }
    & & \\ \\

\end{supertabular}
}

\end{document}
